\chapter{Evaluation}
\label{chap:evaluation}
The method of Chapter~\ref{chap:methodology} is evaluated on four real-life
event logs. The effect of a recommendation on a running case cannot be read
from the historical data, since each case followed only one continuation, so
it is estimated by simulating every case both with and without the
recommendation.


\section{Research Questions}
\label{sec:ev:rq}
% TODO: research questions to be decided with the supervisors.


\section{Case Studies}
\label{sec:ev:cases}
The evaluation uses four event logs, the same as in~\cite{le2025balancing}.
The \emph{Bank Account Closure} (BAC) log records a process of a banking
institution that closes the accounts of its customers, at the request of
either the customer or the bank, and is not publicly available. The
BPIC2012~\cite{bpic2012} and BPIC2017~\cite{bpic2017} logs were published for
the Business Process Intelligence Challenge and record the loan applications
handled by a Dutch financial institution. BPIC2017 contains a concept drift,
so it is split into two logs, \emph{bpi17-before} with the cases before week
22 and \emph{bpi17-after} with the cases after week 28.

Two KPIs are used. The first is the \emph{outcome} of the case, a binary KPI
to be maximized. In BPIC2012, bpi17-before and bpi17-after a case is positive
when the customer accepts an offer, that is, when the activity
\texttt{O\_ACCEPTED} (\texttt{O\_Accepted} in BPIC2017) occurs. In BAC a case
is negative when one of the activities \texttt{Network Adjustment Requested}
and \texttt{Back-Office Adjustment Requested} occurs, since both are
corrective actions that the bank wants to avoid. The second KPI is the
\emph{remaining time} of the case, from the end of the last event of the
prefix to the end of the case, to be minimized. The activities that define
the outcome form the set $\mathcal{A}_F$ of each log and are never
recommended (Section~\ref{subsec:met:offline:feasible}).

The cases with fewer than four events are removed from every log, since they
leave too few prefixes on which to recommend. Only BAC contains such cases,
3\,400 out of 32\,429, and two of its activities, \texttt{Request deleted}
and \texttt{Request completed with customer recovery}, occur only in them,
so the log keeps 13 of its 15 activities.

The events of BPIC2012 and BPIC2017 are of three kinds, distinguished by the
prefix of the activity name: \texttt{A\_} for the states of the application,
\texttt{O\_} for the states of the offers and \texttt{W\_} for the work items
handled by the employees of the bank. All the \texttt{W\_} events are
removed, six activities per log, and only the \texttt{A\_} and \texttt{O\_}
activities are kept. This removes 46\% of the events of BPIC2012 and about
9\% of those of the two BPIC2017 logs. Only rows are dropped: the remaining
events keep their order and attributes, and no case is lost, since none
consists of \texttt{W\_} events alone.
% TODO (user): add the reason why the W_ activities are removed.

In the versions of BPIC2012 and BPIC2017 used here, the start of an event is
not recorded: its start timestamp is a copy of the completion of the previous
event of the case, so every activity seems to begin as soon as the previous
one ends, and no waiting time is visible. The start times are therefore
estimated with the method of Vinci et al.~\cite{vinci2026healthcare}, which
uses only the completion times. The start of an event $e$ of activity $a$ is
placed between its earliest possible start and its completion,
\begin{equation}
  \mathit{start}(e) = \mathit{end}(e) - \alpha(a)\,
  \bigl( \mathit{end}(e) - \mathit{mintime}(e) \bigr),
  \label{eq:ev:start}
\end{equation}
where $\mathit{mintime}(e)$ is the later of the completion of the previous
event of the same case and the completion of the latest earlier event of the
same resource, and $\alpha(a) \in [0, 1]$ is a parameter of the activity:
with $\alpha(a) = 1$ the activity starts as soon as possible, with
$\alpha(a) = 0$ it takes no time (Figure~\ref{fig:ev:start}). The parameters
are searched by bisection over 20 iterations. At every iteration the log is
repaired with the current values, a probability distribution is fitted to the
duration of each activity, and the process is simulated with a model built on
the Petri net with fitness one of Section~\ref{subsec:ev:sim:net}, together
with the resource calendars, the arrival rate and the branching
probabilities discovered from the log. The simulated and the real log are
compared activity by activity on the times between consecutive completions
of a case, through the Wasserstein distance, and each activity keeps the
value of $\alpha(a)$ with the smallest distance. The BAC log records the
start of its events and is left unchanged.

\begin{figure}[h]
  \centering
  \providecolor{Darkcandyapplered}{rgb}{0.64, 0.0, 0.0}
  \begin{tikzpicture}[x=0.9cm,y=0.68cm,font=\small]
    \node[anchor=east] at (0,2) {case of $e$};
    \node[anchor=east] at (0,1) {resource of $e$};
    \draw[->] (0,0.35) -- (12.5,0.35) node[anchor=west] {$t$};
    \draw[fill=gray!25] (0.4,1.75) rectangle (3.4,2.25);
    \node at (1.9,2) {previous event};
    \draw[fill=gray!25] (2.2,0.75) rectangle (5.0,1.25);
    \node at (3.6,1) {previous event};
    \draw[pattern=north east lines, pattern color=Darkcandyapplered!60]
      (5.0,1.75) rectangle (7.8,2.25);
    \draw[fill=Darkcandyapplered!75] (7.8,1.75) rectangle (11.4,2.25);
    \node[white] at (9.6,2) {\textbf{$e$ (estimated)}};
    \draw[dashed] (5.0,0.35) -- (5.0,2.5);
    \draw[dashed] (7.8,0.35) -- (7.8,2.5);
    \draw[dashed] (11.4,0.35) -- (11.4,2.5);
    \node[below] at (5.0,0.35) {$\mathit{mintime}(e)$};
    \node[below] at (7.8,0.35) {$\mathit{start}(e)$};
    \node[below] at (11.4,0.35) {$\mathit{end}(e)$};
    \draw[decorate,decoration={brace,amplitude=4pt}] (5.0,2.6) -- (7.8,2.6)
      node[midway,above=3pt] {waiting};
    \draw[decorate,decoration={brace,amplitude=4pt}] (7.8,2.6) -- (11.4,2.6)
      node[midway,above=3pt,font=\footnotesize]
      {$\alpha(a)\,\bigl(\mathit{end}(e) - \mathit{mintime}(e)\bigr)$};
  \end{tikzpicture}
  \caption{Estimated start of an event $e$ of activity $a$
    (Equation~\eqref{eq:ev:start}). The event cannot start before
    $\mathit{mintime}(e)$, when both the previous event of its case and the
    previous event of its resource have ended. The parameter $\alpha(a)$
    splits the interval up to $\mathit{end}(e)$ into a waiting time and the
    duration of the event.}
  \label{fig:ev:start}
\end{figure}

Table~\ref{tab:ev:logs} summarizes the four logs after these steps.

\begin{table}[h]
  \centering
  \small
  \begin{tabular}{lrrrrrr}
    \toprule
    Log & Cases & Events & Activities & Resources & Mean duration & Positive \\
        &       &        &            &           & (days)        & cases (\%) \\
    \midrule
    BAC          & 29\,029 & 205\,756 & 13 & 654 & 16.2 & 82.4 \\
    BPIC2012     &  4\,685 &  62\,114 & 17 &  60 & 18.6 & 47.9 \\
    bpi17-before & 12\,121 & 166\,547 & 18 & 103 & 20.2 & 56.4 \\
    bpi17-after  & 14\,066 & 192\,723 & 18 & 119 & 21.9 & 54.7 \\
    \bottomrule
  \end{tabular}
  \caption{The four event logs after preprocessing. The duration of a case
    runs from the start of its first event to the end of its last one, with
    the estimated start times for BPIC2012 and the two BPIC2017 logs.}
  \label{tab:ev:logs}
\end{table}

\section{Experimental Setting}
\label{sec:ev:setting}
Each log is split in time. The split time $t_{\mathrm{split}}$ is the
earliest instant at which 80\% of the cases are completed. The cases
completed by $t_{\mathrm{split}}$ form the training log $\mathcal{L}$, from
which the feasibility prescription module and the predictive models are
built. The cases still running at $t_{\mathrm{split}}$ form the test log
$\mathcal{L}^{\mathrm{test}}$, and the cases that start after it are
discarded.

The setting reproduces a user who asks for recommendations at
$t_{\mathrm{split}}$ and knows only what has happened up to then. Every case
of $\mathcal{L}^{\mathrm{test}}$ is therefore cut at $t_{\mathrm{split}}$:
its prefix is made of the events completed by that instant, and the prefixes
form the log $\mathcal{L}^{\mathrm{run}}$ of running cases, on which the
method is applied. The next step of each case is recommended after the last
event of its prefix. The complete cases of $\mathcal{L}^{\mathrm{test}}$,
whose continuation after $t_{\mathrm{split}}$ the method never sees, are used
only to check the simulator (Section~\ref{sec:ev:sim}). A case whose first
event is still running at $t_{\mathrm{split}}$ has no completed event and
therefore no prefix; this happens to 27 cases of BAC, which are left out of
$\mathcal{L}^{\mathrm{run}}$. Table~\ref{tab:ev:split} reports the result of
the split.

\begin{table}[h]
  \centering
  \small
  \begin{tabular}{lcrrr}
    \toprule
    Log & $t_{\mathrm{split}}$ & $\lvert \mathcal{L} \rvert$ &
      $\lvert \mathcal{L}^{\mathrm{run}} \rvert$ & Discarded \\
    \midrule
    BAC          & 2019-01-29 & 23\,224 & 1\,876 & 3\,902 \\
    BPIC2012     & 2012-02-16 &  3\,748 &    623 &    314 \\
    bpi17-before & 2016-05-27 &  9\,697 & 1\,218 & 1\,206 \\
    bpi17-after  & 2016-12-16 & 11\,253 & 1\,840 &    973 \\
    \bottomrule
  \end{tabular}
  \caption{Temporal split of the four logs: number of training cases, of
    running cases with at least one completed event, and of discarded cases,
    which start after $t_{\mathrm{split}}$.}
  \label{tab:ev:split}
\end{table}


\section{Offline Phase}
\label{sec:ev:offline}

\subsection{Feasibility Prescription Module}
\label{subsec:ev:off:feasible}
The transition system is built on $\mathcal{L}$ with a window of $w = 5$
activities. The module is then restricted further than in
Section~\ref{subsec:met:offline:feasible}, because a step that occurs only a
handful of times in $\mathcal{L}$ is scored by the predictive models with
almost no data behind it, and its predictions are unreliable. A next
activity is kept only if it follows the window of the prefix in at least 1\%
of the occurrences of that window in $\mathcal{L}$. A resource is kept for a
next activity only if it performs that activity after the same window in at
least 5\% of its occurrences there, and the most frequent resource is always
kept, so that this second threshold never removes an activity. As a
consequence, the eligible resources of an activity depend on the window of
the prefix, and not only on the activity as in
Equation~\eqref{eq:met:off:resources}. In BAC, 7\% of the training events
carry the resource value \texttt{NotDef}, which marks an event without a
defined resource; it is treated like any other resource.

\subsection{Predictive Models}
\label{subsec:ev:off:models}
Both predictive models are gradient-boosted trees built with
CatBoost~\cite{prokhorenkova2018}, which runs in its Plain mode. The outcome
is learned by a classifier trained with the logarithmic loss, the remaining
time by a regressor trained with \texttt{RMSEWithUncertainty}, which predicts
a mean and a variance (Section~\ref{subsec:bg:gbm:catboost}). Both are
trained with stochastic gradient Langevin boosting and used as virtual
ensembles (Section~\ref{subsec:bg:uncertainty:gbm}). Each requests $M = 50$
members, and fewer when the model has fewer than 100 trees, since a member
needs at least two of them. The input is the vector of
Section~\ref{subsec:met:offline:training}: the categorical features, among
them the activity and resource of the last event, the next activity and
resource and the case attributes of each log, are one-hot encoded, and the
continuous ones are standardized. The case attributes are the requested
amount in BPIC2012; the requested amount, the loan goal and the application
type in the two BPIC2017 logs, together with the event attributes
\texttt{Action} and \texttt{EventOrigin}; the closure type and reason and the
role of the resource in BAC.

The regressor does not model the remaining time in seconds. Remaining times
are skewed, with a long tail of slow cases, so the Normal distribution of
\texttt{RMSEWithUncertainty} is fitted to
\begin{equation}
  y = \log\bigl( 1 + \text{remaining time in days} \bigr).
  \label{eq:ev:timespace}
\end{equation}
The predictive distribution of the remaining time is then log-normal. It is
truncated at the longest remaining time in $\mathcal{L}$, because a large
predicted variance would otherwise inflate its mean $e^{\mu + \sigma^2/2}$
without bound, and the point prediction is the mean of the truncated
distribution, converted back to seconds (Appendix~\ref{app:lognormal}). The confidence filter compares the
members directly in the space of~\eqref{eq:ev:timespace}: the transformation
is monotone, so a candidate that is faster than the baseline there is faster
in seconds too.

The hyperparameters of each model are chosen with Optuna~\cite{akiba2019}
over 80 trials, in the search space of Table~\ref{tab:ev:search}. A trial
trains a model on the training part of
$\mathcal{D} \cup \mathcal{D}^{\varnothing}$ and measures its loss on the
validation set, which holds the last 20\% of the cases of $\mathcal{L}$ by
completion time, with all their prefixes and their copies in
$\mathcal{D}^{\varnothing}$. A trial grows at most 3\,000 trees and stops
early when the validation loss has not improved for 50 of them, and a trial
whose loss after 100 trees is worse than the median of the earlier trials
at the same point is abandoned. The trial with the lowest validation loss
wins. The final model is then refitted on all of
$\mathcal{D} \cup \mathcal{D}^{\varnothing}$ with the winning
hyperparameters and with the number of trees at which the winning trial
stopped.

\begin{table}[h]
  \centering
  \small
  \begin{tabular}{lll}
    \toprule
    Hyperparameter & Range & Notes \\
    \midrule
    learning rate $\epsilon$ & $[0.02, 0.3]$, log scale & \\
    tree depth & $\{3, \dots, 12\}$ & \\
    L2 coefficient $\zeta$ & $[1, 50]$, log scale & \\
    share of features per split & $[0.6, 1]$ (classifier), $[0.5, 1]$ (regressor) & \\
    sampling of the examples & Bayesian, Bernoulli, MVS & \\
    \quad bagging temperature $\omega$ & $[0, 10]$ & Bayesian only \\
    \quad sampling rate $q$ & $[0.5, 1]$ & Bernoulli and MVS \\
    growing policy & oblivious, depthwise & \\
    \quad minimum examples per leaf & $\{1, \dots, 200\}$, log scale & depthwise only \\
    \bottomrule
  \end{tabular}
  \caption{Search space of the hyperparameters, the same for the four logs.
    The symbols are those of Section~\ref{subsec:bg:gbm:catboost}.}
  \label{tab:ev:search}
\end{table}

The temperature $T$ of the classifier and the factor $s$ of the regressor
(Section~\ref{subsec:met:offline:models}) are fitted on the validation set
with the model of the winning trial and then applied to the final model.
This assumes that refitting on slightly more data does not change how
miscalibrated the model is.

The models are assessed on the prefixes of the cases in
$\mathcal{L}^{\mathrm{test}}$, which none of them has seen. For the
classifier the measures are the logarithmic loss and the expected
calibration error over ten equal-width bins
(Equation~\eqref{eq:bg:unc:ece}), before and after temperature scaling,
together with the average data and knowledge uncertainty of its virtual
ensemble. For the regressor they are the mean absolute error and the root
mean squared error of the point prediction, in days; the ENCE over fifteen
bins of equal size (Equation~\eqref{eq:bg:unc:ence}), the coefficient of
variation of the predicted standard deviations and the coverage of the
$\pm\sigma$ and $\pm 2\sigma$ intervals, before and after the scaling by
$s$; and the average data and knowledge uncertainty of its ensemble. The
calibration measures of the regressor are computed in the space
of~\eqref{eq:ev:timespace}, where the predictive distribution is Normal.


\section{Online Phase}
\label{sec:ev:online}
For every prefix of $\mathcal{L}^{\mathrm{run}}$, all the candidates given
by the feasibility prescription module are evaluated, together with the
baseline. The outcome is to be maximized and the remaining time to be
minimized. The confidence filter uses $\gamma_j = 0.5$ for both KPIs, so a
candidate passes when, for at least one of them, an improvement over the
baseline is more likely than not. The selection rule then returns up to
$k = 5$ candidates of the Pareto front, or the whole front when it has fewer
points. The candidates of a case are a set: the order in which they are
listed carries no preference, and each of them is evaluated on its own.


\section{Simulation}
\label{sec:ev:sim}
Whether a recommendation helps could be measured directly only by an A/B
test, in which some running cases follow it and others do not, and this is
not possible on historical data. Every prefix of $\mathcal{L}^{\mathrm{run}}$
is instead continued by a simulator, once without a recommendation and once
for each of its recommended steps, and the continuations are compared. The
simulator plays no part in producing the recommendations: it is used only
here.

\subsection{Petri Net}
\label{subsec:ev:sim:net}
The simulator needs a Petri net that can replay every prefix and on which
every recommended step can be taken. A net with fitness one on the whole log
meets both requirements. It is built from the directly-follows graph of the
log, the graph with one node per activity and an arc from $a$ to $b$
whenever $b$ immediately follows $a$ in some case. Each activity becomes a
place, meaning that $a$ was the last activity executed, and each arc
becomes a transition labeled $b$ that moves the token from the place of $a$
to the place of $b$. A transition labeled with each first activity of the
cases leaves the initial place, and a silent transition leads from the place
of each last activity to the final place. Every case of the log is then a firing sequence of the net, so its
fitness is one. Every recommended step can be taken as well: the next
activities of a prefix ending with $a$ are those observed after a window
ending with $a$ (Section~\ref{subsec:met:offline:feasible}), so the
corresponding transitions are enabled in the place of $a$. The price is a
lower precision, since the net remembers only the last activity of a case:
it cannot represent concurrency, and it allows sequences that respect every
observed pair of consecutive activities but never occurred as a whole.

The net is built from the whole log, with the preprocessing of
Section~\ref{sec:ev:cases} except the removal of the short cases, and so are
the parameters of the simulator. The simulator therefore sees the test cases
too, but only through aggregate statistics, and in the same way for the
runs without and with a recommendation.

\subsection{Simulation Model}
\label{subsec:ev:sim:model}
The simulation model is discovered with an adapted version of
ProSiT~\cite{vinci2025prosit}. Waiting and routing are modeled place by
place, time first. For every place a regression tree predicts, on the scale
$\log(1 + \text{days})$, how long a case waits after its last event before
its next one starts, from the attributes of the case, the activities it has
executed so far and its age. For every place with more than one outgoing
transition, a classification tree then chooses the transition from the same
information and from the wait just drawn, so that what happens next can
depend on how long the case has waited. The execution time of each activity
is predicted by a regression tree of depth up to three, chosen by
cross-validation. In every tree, a new value is drawn among the values
observed in the training examples of the leaf the case falls into, rather
than from a distribution fitted to them. The resource of an activity is
drawn with the frequencies observed in the log, and every resource follows
the working calendar discovered for it.

\subsection{Baseline and Recommendation Runs}
\label{subsec:ev:sim:runs}
For every case of $\mathcal{L}^{\mathrm{run}}$ the simulator replays the
prefix on the net and continues the case from the marking it reaches. The
continuation starts right after the last event of the prefix, with no
constraint tied to $t_{\mathrm{split}}$, so the simulator starts from exactly
the information available to the method. No new case arrives during the
simulation: only the cases of $\mathcal{L}^{\mathrm{run}}$ are continued.

In the \emph{baseline} runs the case continues freely until it completes. In
a \emph{recommendation} run its next event is the recommended step: the
recommended activity, performed by the recommended resource. The wait before
it is drawn from the waits observed in the current place, weighted by how
likely the recommended transition is after each of them, and the rest of the
case continues freely. A recommendation is applied only if its activity is
the next visible step of the case, enabled in the current marking or
reachable through silent transitions alone; a run in which it is not
applicable is left out of the evaluation. Each prefix is simulated ten times
without a recommendation and ten times for each of its recommended steps.
From every run the two KPIs are read as in Section~\ref{sec:ev:cases}: the
outcome of the simulated case and its remaining time, measured from the end
of the last event of the prefix, the same reference as the predictions.

\subsection{Measures}
\label{subsec:ev:sim:measures}
A recommended step of a case is a \emph{point}. For each point the effect of
the recommendation is measured twice, by the predictive models and by the
simulation, and in both cases against the baseline. For the outcome the
effect is the change of the probability of a positive outcome, in
percentage points: on the predicted side the difference between the
probabilities predicted for $x_C$ and for $x_B$, on the simulated side the
difference between the shares of positive runs with the recommendation and
without it. For the remaining time it is the relative change, in percent of
the baseline: on the predicted side between the predictions for $x_C$ and
$x_B$, on the simulated side between the mean remaining times of the two
groups of runs.

A point \emph{improves} a KPI when its change is strictly favorable, a
higher outcome or a shorter remaining time. On the simulated side, a case
whose runs are all positive both with and without the recommendation counts
as improved on the outcome, since it cannot do better. Comparing the two
sides gives, for each KPI and for the two KPIs together, a two-by-two table:
among the points that the models predict to improve, the share that also
improves in the simulation, and the same share among the points that the
models predict not to improve. If the prediction carried no information the
two shares would be equal, so their difference measures how informative the
prediction is.

A point is \emph{good} when it improves one KPI by at least $x$ and worsens
the other by at most $y$, with $x = 5$ and $y = 10$, in percentage points
for the outcome and in percent for the remaining time; the same rule is
applied to the predicted and to the simulated changes. Finally, the expected
value of each KPI with the recommendation is computed on both sides,
averaging first over the points of a prefix and then over the prefixes. It
is computed on all the points and on the good points alone, as if only those
were recommended, together with the share of prefixes left with at least one
good point.

\subsection{Checking the Simulator}
\label{subsec:ev:sim:check}
The comparison above is meaningful only if the simulator reproduces the
process. The baseline runs are therefore compared with what actually
happened to the same cases, the complete cases of
$\mathcal{L}^{\mathrm{test}}$: the share of positive cases and the mean,
median and standard deviation of their duration, also broken down by the age
of the case at $t_{\mathrm{split}}$.

\subsection{A Fully Simulated Training Log}
\label{subsec:ev:sim:full}
When the predicted and the simulated effects of a recommendation disagree,
the cause may lie in the predictive models or in the simulator, which may
not reproduce the real process. To tell the two apart, the evaluation is
repeated on BPIC2012 with a training log produced by the simulator itself.
The simulation model discovered from BPIC2012 generates as many cases as its
training log contains, 3\,748, each from the initial marking of the net to
its end, with arrival times drawn from the discovered arrival process and
case attributes drawn from their distribution in the log; no real prefix is
replayed. The resulting log, \emph{BPI12-sim}, takes the place of the
training log, while $\mathcal{L}^{\mathrm{test}}$ and
$\mathcal{L}^{\mathrm{run}}$ stay those of the real log, so the method is
asked for recommendations on the same running cases. The whole method is run
again on BPI12-sim: the feasibility prescription module and the predictive
models are built on the simulated training log, and the recommendations are
simulated with a Petri net and simulator parameters discovered from the
simulated log. In this setting the predictive models and the simulator learn
the same process, the one encoded in the simulation model, so a disagreement
between them can no longer come from a simulator that departs from the real
process. The simulated log is also compared with the real BPIC2012 log, with
the measures of Section~\ref{subsec:ev:sim:check}.
